\documentclass[11pt]{article}
\usepackage[margin=.85in]{geometry}
\usepackage[T1]{fontenc}
\usepackage{amsmath,amssymb,amsthm,booktabs,array,hyperref,microtype,seqsplit}
\hypersetup{colorlinks=true,urlcolor=blue,linkcolor=blue,citecolor=blue}
\setlength{\emergencystretch}{2em}
\newtheorem{theorem}{Theorem}
\newtheorem{proposition}[theorem]{Proposition}
\newtheorem{lemma}[theorem]{Lemma}
\newcommand{\Ftwo}{\mathbb F_2}
\newcommand{\poly}{\operatorname{poly}}
\title{LSPV-4: PAID-SELECTOR-GADGET-FIREWALL224\\[3pt]
\large Connected 3SAT as Disjoint-XOR--2CNF With a Linear Projected-Excess Parameter}
\author{LSPV-4, Independent Time--Space and Circuit Resource Auditor}
\date{29 October 2026, 07:36 PDT}
\begin{document}
\maketitle
\begin{abstract}
We independently audit the proposed adaptive $\min(k,K)$ mixed-affine/2SAT portfolio. An explicit uniformly computable gadget maps any proper signed 3CNF $F$ to disjoint ternary affine equations plus ordinary signed 2CNF, preserving the \emph{complete} Boolean function after existentially forgetting the selectors. For every connected $m$-clause source, the resulting mixed system has exactly one affine--binary interaction cluster, projected affine excess $K=m$, and binary-graph vertex-cover number $2m\le k\le3m$. Consequently, the previously proved portfolio upper bound is $2^m\poly(N')$ on this image, including for trivial all-positive SAT sources. The image class is NP-complete. These are exact \emph{structural parameter} statements and a hardness-preserving reduction, not time lower bounds on all solvers. We give distinct single-tape, two-tape, RAM and efficient fixed-universal-machine ledgers, reproduce 69,904 elimination-control assignments, and independently rerun the Library-connected-interface audit. $P$ versus $NP$ remains OPEN.
\end{abstract}

\section{PROPOSE: falsifiable missing mechanisms}
The current gap is not merely a missing Gaussian or SCC implementation. It is the absence of an unconditional, all-input transformation of 3SAT into a representation with small \emph{charged} adaptive parameters and a strictly decreasing complete encoded size. We test four mechanisms. \textbf{M1}: a generic polynomial 3SAT-to-affine--2CNF compiler might automatically give small projected excess; the gadget below falsifies this. \textbf{M2}: a small binary vertex cover could rescue the same compiler; disjoint selector triangles force $k\ge2m$. \textbf{M3}: a high measured $\min(k,K)$ might certify hardness even on a connected family; positive-chain controls falsify this inference because an all-ones assignment decides them immediately. \textbf{M4}: the elementary static connected-interface obstruction survives on its exact syntactic model but says nothing about the adaptive mixed compiler's run time; direct proof and rerun confirm its scope.

\section{SEARCH: actual source and applicable hypotheses}
The exact source for the predecessor is 
\href{https://github.com/j15marti/pnp/blob/main/research/agents/lspv-4/2026-10-29_0636_LSPV-4_ADAPTIVE_MIN_CLUSTER_PORTFOLIO223.tex}{LSPV-4 predecessor report}, repository \texttt{j15marti/pnp}, path \path{research/agents/lspv-4/2026-10-29_0636_LSPV-4_ADAPTIVE_MIN_CLUSTER_PORTFOLIO223.tex}, blob 
\texttt{\seqsplit{04d9a94e070c93a37e29d079cbd0676d2eda1203}}. Its theorem concerns mixed affine equations and binary clauses with components formed from XOR support, then clusters formed by binary interactions. It gives a \emph{uniform upper bound} $\poly(N)(1+\sum_j2^{\min(k_j,K_j)})$ with vertex-cover discovery charged. We do not turn that upper bound into a lower bound.

The actual \href{https://github.com/msoos/cryptominisat/blob/master/src/gaussian.cpp}{CryptoMiniSat \path{src/gaussian.cpp}}, blob \texttt{\seqsplit{7b9d200c1c1d8e344ce0c3109bfc92806a14359b}}, is an implementation of mixed Boolean/XOR engineering, not a proof of a polynomial solver on arbitrary mixed constraints. The actual \href{https://github.com/uds-psl/coq-library-complexity/blob/coq-8.16/theories/NP/TM/mTM_to_singleTapeTM.v}{Coq library \path{theories/NP/TM/mTM_to_singleTapeTM.v}}, blob \texttt{\seqsplit{a12c882a5aed008870e913dd292b6f2530432ecc}}, exhibits explicit tape encoding and polynomial reductions; it does not certify the exponents below and no kernel compilation was attempted \cite{gahekunze}. The actual \href{https://github.com/openai/math/blob/main/preprints/Polynomial-Time-Unitary-Synthesis-from-a-Boolean-Oracle-October-5-2026/build/sections/01-introduction.tex}{\texttt{openai/math/.../01-introduction.tex}}, blob \texttt{\seqsplit{5ad06c26238e2572e93f68b380f542b2ab1845fe}}, explicitly places no polynomial classical circuit-size bound on the unitary-dependent oracle and does not construct a SAT shrinker. GitHub searches included ``exactly one pairwise at most one xor parity SAT'', ``XOR clause CNF two sat gaussian'', ``CardEnc.atmost1'', and ``XOR SAT 2SAT combined NP complete''. We compared the construction against Schaefer's Boolean constraint dichotomy \cite{schaefer} and heterogeneous backdoor literature \cite{gaspers}. Repository matches and code tests are not kernel-certified asymptotic theorems.

\section{BUILD: exact projection and parameter inflation}
Let $F(x)=\bigwedge_{j=1}^{m}(\ell_{j1}\lor\ell_{j2}\lor\ell_{j3})$ be proper signed 3CNF, with three distinct variable names per clause and no unused variables. Introduce fresh, pairwise distinct selectors $s_{j1},s_{j2},s_{j3}$ for each clause. Put
\begin{align}
E_j&:\quad s_{j1}\oplus s_{j2}\oplus s_{j3}=1,\\
B_j&:\quad\bigwedge_{1\le a<b\le3}(\neg s_{ja}\lor\neg s_{jb})
\;\land\;\bigwedge_{a=1}^{3}(\neg s_{ja}\lor\ell_{ja}),\\
G_F(x,s)&:=\bigwedge_{j=1}^{m}(E_j\land B_j).\label{eq:gadget}
\end{align}
Every affine equation has three distinct variables and distinct affine equations have disjoint variable supports. Every ordinary Boolean clause in $B_j$ has width exactly two and distinct variable names. There are exactly $m$ ternary parity equations, $6m$ binary clauses and $n+3m$ variables.

\begin{theorem}[Exact projected semantics, PROVED/RESTRICTED]\label{th:projection}
For every assignment $x\in\{0,1\}^n$,
\[
\boxed{\quad (\exists s\;G_F(x,s))\ \Longleftrightarrow\ F(x).\quad}
\]
The transformation is deterministic, uniform, polynomial time, and uses no SAT oracle or answer-side bit.
\end{theorem}
\begin{proof}
An odd-parity triple together with all three negative pairwise clauses admits exactly the three assignments $100,010,001$: parity excludes $000$, while the pairwise clauses exclude any assignment with two or three ones. If the $a$th selector is one, its implication $(\neg s_{ja}\lor\ell_{ja})$ forces $\ell_{ja}$ true. Hence a selector extension exists for clause $j$ only if that original clause is satisfied. Conversely, for a satisfied original clause choose one true literal and set exactly its selector to one; this satisfies all seven constraints of the local gadget (one parity equation and six binary clauses). The selector sets for different $j$ are disjoint. Therefore all gadget blocks extend simultaneously if and only if all source clauses are satisfied.
\end{proof}

Let the \emph{affine components} be connected components of XOR support, treating each original $x_i$ not occurring in an XOR as a singleton. For component $C_i$, let $P_i$ be all its variables appearing in a binary clause, let $L_i$ span its homogeneous affine solution space, and write $r_i=\operatorname{rank}_{\Ftwo}(\pi_{P_i}L_i)$. As in the predecessor, set $K=\sum_i\max(r_i-1,0)$ over a binary-interaction cluster. Let $k$ be minimum vertex-cover size of the \emph{binary variable graph} in that cluster.

\begin{theorem}[One connected cluster with linear parameter, PROVED/RESTRICTED]\label{th:params}
If the incidence graph of $F$ is connected, the affine--binary interaction graph of $G_F$ is connected. Its unique cluster satisfies
\[
\boxed{\ K=m,\qquad2m\le k\le3m,\qquad\min(k,K)=m.\ }
\]
Thus the specific adaptive $\min(k,K)$ portfolio proves only the upper bound $2^m\poly(N_G)$ on this image. This expression is not a lower bound on its execution time, and is not a lower bound on arbitrary SAT algorithms.
\end{theorem}
\begin{proof}
For each $j$, the XOR support is a triple with a single independent parity equation, hence two free affine degrees of freedom. All three selectors occur in the negative-pairwise clauses, so $r_j=2$ and the contribution is one. Each original variable belongs to a singleton affine component, appears in an implication, has projected rank one, and contributes zero. Summing gives $K=m$.

The graph of binary clauses contains one triangle on $s_{j1},s_{j2},s_{j3}$ for each $j$, and the selector triangles are vertex-disjoint. Each triangle needs at least two vertices in any vertex cover, so $k\ge2m$. Selecting all $3m$ selectors covers both the triangle edges and the implication edges, hence $k\le3m$. Contract each selector triple to its affine component. The implication edges connect this component to every original variable of clause $j$. Thus the contracted interaction graph is exactly a subdivision of the source incidence graph, which is connected. Therefore there is one cluster and $\min(k,K)=m$.
\end{proof}

\begin{proposition}[NP-complete restricted image, PROVED/RESTRICTED]
Deciding satisfiability of systems having the exact pattern \eqref{eq:gadget} is NP-complete, even with connected source incidence, despite the disjoint ternary affine equations. The proof is a polynomial reduction from ordinary proper 3SAT by Theorem~\ref{th:projection}, and membership in NP follows from evaluation of an explicitly given $(x,s)$ assignment. To enforce connectivity without deciding $F$, first remove unused source variables and connect the incidence components: add a common new variable $z$ and a fresh $w_i$ for each original component, then add a proper clause $(z\lor v_i\lor w_i)$ containing a chosen old variable $v_i$ from that component. Setting $z=1$ satisfies all bridges, so satisfiability is preserved. The result is connected. This construction does not retain strict signed exact-$(2,2)$ occurrence counts.
\end{proposition}

\paragraph{A positive-family falsifier for an invalid separation inference.}
Set $F_m:=\bigwedge_{j=1}^{m}(x_j\lor x_{j+1}\lor x_{j+2})$ on $m+2$ variables. This is proper, connected, and visibly satisfiable by the all-ones assignment. Nevertheless its mixed image has $K=m$ and $k\ge2m$. Large structural portfolio parameters therefore cannot be interpreted as an algorithm-independent hard-instance certificate.

\section{Full encoded-length and native machine ledgers}
Use an unpadded canonical source containing a header, $m$ explicitly signed triple clauses and variable identifiers. For the mixed child, print a header, the $m$ XOR equations, all $6m$ binary clauses, and identifiers up to $n+3m$. The output length obeys
\[
N_G=\Theta((n+m)\log(n+m+2))\quad\text{for families with }m=\Omega(n),
\qquad N_G=O(N\log(N+2))
\]
for general well-formed source encodings with no unused variables and explicit clauses. In the positive chain, $N_G=\Theta(N)$ but the observed byte constant is greater than one. This is \emph{not} a shrinking transformation. Its output grammar also differs from its input grammar, so it fails the required recursive closure condition.

Let $M=N+N_G+2$. A deliberately conservative sequential compiler repeatedly scans the source, emits each equation and binary clause, and keeps binary counters and addresses. Its fixed read-only-input multitape bound is $T_k=O(M^2)$ and $S_k=O(N_G+\log M)$ writable cells, counting the materialized output. For two fully writable tapes, the same simple scanning implementation has $T_2=O(M^2)$, $S_2=O(N_G+\log M)$ auxiliary visited cells and $D_{\mathrm{total}}\le T_2$ symbol changes; no smaller write cap was certified. With the stipulated conservative one-tape simulation,
\[
T_1=O(T_k(N+S_k+\log(T_k+2)))=O(M^3),\qquad
P_1=O(N+S_k+\log(T_k+2))=O(M),
\]
with original input included in $P_1$. An indexed word RAM with $\Theta(\log M)$-bit words can perform source parsing and explicit output in $O(M\operatorname{polylog} M)$ word operations, but this is a separate model and not a RAM-native SAT lower bound. A fixed universal TM selected to have polynomial simulation overhead executes in $O(M^C)$ for some compiler- and universal-machine-dependent constant $C$; an arbitrary inefficient universal interpreter need not satisfy that claim. No precise universal exponent or Coq kernel certificate is asserted.

One fixed converter handles both answer classes with the same bounds. Denote its restrictions on satisfiable and unsatisfiable sources by $C_+$ and $C_-$. Both have the same \emph{upper} envelope $O(M^2)$ on the fixed-multitape construction; neither receives the answer side for free. The successor mixed formula has the same decision on both sides but, under the predecessor portfolio, only the \emph{guaranteed upper bound} $2^m\poly(N_G)$ applies to either. Side-specific worst cases and same-input simultaneous time-space extrema are not established. None of the tape transport restrictions is used to infer $P\ne NP$.

\section{Independent computational and source audit}
The standalone Python reproducer, seed 20261029, completed 5,488 all-original-assignment clause-extension checks, 135,424 independently evaluated full mixed-assignment checks, 2,793 connectedization witnesses, three exhaustive small-graph vertex-cover nonexistence checks, and 32 matched SAT/UNSAT cofactor comparisons, all passing. Matched connected proper 3CNF inputs with $n=4,m=9$ have equal physical source length (124 bytes) and equal mixed-child length (713 bytes), with different SAT answers and the same $K=9$, $k\ge18$. Positive-chain examples through $m=192$ were independently checked for one cluster, $K=m$, the explicit cover bounds, and actual printed byte sizes.

We also independently reproduced the classical single-variable elimination identity. With $A_i=(a_i\lor b_i)$ and $B_j=(c_j\lor d_j)$, all surviving variables disjoint,
\[
\exists z\left[\bigwedge_i(z\lor A_i)\land\bigwedge_j(\neg z\lor B_j)\right]
=\bigwedge_{i,j}(A_i\lor B_j)
=\left(\bigwedge_iA_i\right)\lor\left(\bigwedge_jB_j\right).
\]
For $t=1,2,3,4$, all $16+256+4096+65536=69{,}904$ assignments to the surviving variables matched all three expressions. The $t^2$ resolvents versus $O(t)$ circuit example proves no inherent projection expansion or general SAT contraction.

The actual Library file \texttt{Connected\_Interface\_Audit.zip}, library file ID \texttt{\seqsplit{libfile\_a2862f0165b4819195a88d9e896a42a7}}, was materialized. Its \texttt{REPORT.tex}, rendered five-page PDF, \texttt{audit.py} and \texttt{results.json} were inspected; reexecuting the original script produced byte-identical JSON, SHA-256 \texttt{\seqsplit{05675f876d6624073507df1e9cedd18a188143db687e454bc2942143d36cbaca}}. This reproduces 4,680 boundary subsets, 7,191 local-discovery checks and 139 brute-force comparisons. Re-derivation of its exact static lemma gives $k\le4b$ and $n\le3b+12bq$; maximizing $3b+12bq$ for fixed $b+q$ yields $b+q\ge\sqrt{n/3}-1/4$. Its static syntactic scope is retained. These finite checks do not imply a general runtime lower bound.

\paragraph{Status and next role-specific obligation.}
\textbf{PROVED/RESTRICTED}: the gadget's exact projection, $K=m$, $2m\le k\le3m$, connectivity, and NP-complete image via a textbook reduction. \textbf{COMPUTATIONAL}: all finite tests and the original Library audit rerun passed. \textbf{CONDITIONAL}: explicit machine simulation exponents hold for the stated compiler model; no Coq kernel proof. \textbf{FALSIFIED}: the assertion that a generic move into a mixed affine--2CNF grammar necessarily produces low $K$ or low vertex cover, or that high $K$ proves algorithm-independent hardness. \textbf{OPEN}: a uniform polynomial-time, answer-preserving strict net contraction or polynomial-value rank descent on arbitrary 3SAT, and any unrestricted separation. The next LSPV-4 target is to audit a proposed \emph{parameter-reducing} transformation on this exact hard image, charging all source access, failed choices, intermediate storage, output materialization, and reconstruction on one fixed SAT/UNSAT program. $P=NP$ and $P\ne NP$ remain open.

\begin{thebibliography}{9}
\bibitem{schaefer} T. J. Schaefer, ``The Complexity of Satisfiability Problems,'' in \emph{Proceedings of the Tenth ACM Symposium on Theory of Computing}, 1978, pp. 216--226. DOI: 10.1145/800133.804350.
\bibitem{gaspers} S. Gaspers, N. Misra, S. Ordyniak, S. Szeider, and S. Zivny, ``Backdoors into Heterogeneous Classes of SAT and CSP,'' \emph{AAAI}, 2014. DOI: 10.1609/aaai.v28i1.9111.
\bibitem{gahekunze} L. G\"aher and F. Kunze, ``Mechanising Complexity Theory: The Cook-Levin Theorem in Coq,'' \emph{ITP}, 2021. DOI: 10.4230/LIPIcs.ITP.2021.20.
\end{thebibliography}
\end{document}